% markpage — export LaTeX (cible : xelatex ou lualatex).
% La compilation avec pdflatex échouera : fontspec + UTF-8 natif
% dans les blocs de code requièrent xelatex / lualatex.
\documentclass[11pt,a4paper]{article}
\usepackage{fontspec}
\IfFontExistsTF{DejaVu Serif}{\setmainfont{DejaVu Serif}}{}
\IfFontExistsTF{DejaVu Sans}{\setsansfont{DejaVu Sans}}{}
\IfFontExistsTF{DejaVu Sans Mono}{\setmonofont{DejaVu Sans Mono}}{}

\IfFileExists{french.ldf}{\usepackage[french]{babel}}{}
\usepackage{amsmath,amssymb,amsthm}
\usepackage{stmaryrd}
\usepackage{graphicx}
\usepackage[export]{adjustbox}
\usepackage{hyperref}
\usepackage{xcolor}
\usepackage{listings}
\usepackage{enumitem}
\usepackage{booktabs}
\usepackage{tabularx}
\usepackage[normalem]{ulem}
\usepackage[breakable,skins]{tcolorbox}
\usepackage{newunicodechar}

% DejaVu Serif (our main text font) doesn't ship the Unicode
% math-bracket / logic glyphs that show up in our prose (e.g.
% explanations of the editor's ligatures). Redirect each known
% glyph to its LaTeX command via \ensuremath, so it renders
% from the math font even when typed in normal text.
\newunicodechar{⟦}{\ensuremath{\llbracket}}
\newunicodechar{⟧}{\ensuremath{\rrbracket}}
\newunicodechar{⟨}{\ensuremath{\langle}}
\newunicodechar{⟩}{\ensuremath{\rangle}}
\newunicodechar{⊢}{\ensuremath{\vdash}}
\newunicodechar{⊣}{\ensuremath{\dashv}}
\newunicodechar{⊨}{\ensuremath{\models}}
\newunicodechar{⊥}{\ensuremath{\bot}}
\newunicodechar{⊤}{\ensuremath{\top}}
\newunicodechar{∀}{\ensuremath{\forall}}
\newunicodechar{∃}{\ensuremath{\exists}}
\newunicodechar{∈}{\ensuremath{\in}}
\newunicodechar{∉}{\ensuremath{\notin}}
\newunicodechar{∅}{\ensuremath{\emptyset}}
\newunicodechar{⊂}{\ensuremath{\subset}}
\newunicodechar{⊆}{\ensuremath{\subseteq}}
\newunicodechar{⊃}{\ensuremath{\supset}}
\newunicodechar{⊇}{\ensuremath{\supseteq}}
\newunicodechar{∪}{\ensuremath{\cup}}
\newunicodechar{∩}{\ensuremath{\cap}}
\newunicodechar{∧}{\ensuremath{\land}}
\newunicodechar{∨}{\ensuremath{\lor}}
\newunicodechar{¬}{\ensuremath{\neg}}
\newunicodechar{★}{\ensuremath{\bigstar}}
\newunicodechar{✓}{\ensuremath{\checkmark}}
\newunicodechar{✗}{\ensuremath{\times}}
\newunicodechar{♥}{\ensuremath{\heartsuit}}
\newunicodechar{♦}{\ensuremath{\diamondsuit}}
\newunicodechar{♠}{\ensuremath{\spadesuit}}
\newunicodechar{♣}{\ensuremath{\clubsuit}}
\newunicodechar{⚠}{\textbf{[!]}}

\hypersetup{colorlinks=true, linkcolor=blue!50!black, urlcolor=blue!50!black}
\lstset{
  basicstyle=\ttfamily\small,
  breaklines=true,
  frame=single,
  framesep=4pt,
}

% Theorem-like environments matching the ::: admonitions (§17 of
% the markpage SPEC). The shared counter `theorem` keeps lemma /
% proposition / corollary numbered in the same sequence as the
% theorems themselves, which is the usual mathematical convention.
\newtheorem{theorem}{Théorème}
\newtheorem{lemma}[theorem]{Lemme}
\newtheorem{proposition}[theorem]{Proposition}
\newtheorem{corollary}[theorem]{Corollaire}
\theoremstyle{definition}
\newtheorem{definition}{Définition}
\newtheorem{example}{Exemple}
\newtheorem{remark}{Remarque}

\title{bda — Block-Diagram Algebra (à la Faust)}
\author{Test Author \\ Test Org}
\date{2026-01-01}

\begin{document}
\maketitle

Le fence \texttt{```bda} accepte une expression construite avec les cinq
opérateurs de composition de l'algèbre de Faust : \texttt{,} parallèle, \texttt{:}
séquentiel, \texttt{<:} split, \texttt{:>} merge, et \texttt{\textasciitilde{}} récursion. Les primitives
sont des identifiants (avec arité par défaut \texttt{(1,1)}, ou annotée
\texttt{Label[n,m]}), des nombres \texttt{(0,1)}, les opérateurs \texttt{+ - * /} \texttt{(2,1)},
les fonctions \texttt{sin cos exp …} \texttt{(1,1)}, ainsi que \texttt{\_} (identité \texttt{(1,1)})
et \texttt{!} (cut \texttt{(1,0)}).


\section{Accumulateur (compteur)}

L'exemple canonique de l'algèbre : un nombre constant ré-injecté dans
un \texttt{+} rebouclé sur lui-même via \texttt{\textasciitilde{} \_}.


\begin{lstlisting}
1 : +~_
\end{lstlisting}


Voir \textbackslash{}ref\{fig:acc\}.


\section{Accumulateur 2 — récursion avec B composé}

Vérifie que les labels du bloc droit (mis en miroir au-dessus) restent
lisibles malgré la réflexion horizontale.


\begin{lstlisting}
(_,1 : +)~(sin : cos : tan)
\end{lstlisting}


\section{Accumulateur — style Faust avec marqueurs \texttt{z⁻¹}}

Mêmes circuits, mais l'option \texttt{delays} (alias \texttt{faust}) place un petit
carré blanc au milieu de chaque fil de feedback B→A pour matérialiser le
délai unitaire implicite du \texttt{\textasciitilde{}}.


\begin{lstlisting}
1 : +~_
\end{lstlisting}


\begin{lstlisting}
A[3,3] ~ B[2,2]
\end{lstlisting}


\section{Passthrough — chaînes de \texttt{\_} en seq}

Vérifie l'optimisation qui shunte un bloc \texttt{\_,\_,\_,…} quand il sert juste
de bundle d'identité dans un seq. Sinon chaque \texttt{\_} génère un coude à
gauche et un coude à droite, créant des escaliers visuels.


\begin{lstlisting}
A[2,3], B[2,3] : _,_,_,_,_,_ : C[6,1]
\end{lstlisting}


\section{Double récursion}

Cas piégeux du wrap des labels : avec \texttt{\textasciitilde{}} à droite-assoc, ceci parse
en \texttt{(\_,1:+) \textasciitilde{} (+ \textasciitilde{} (sin:cos:tan))}. Le sous-bloc \texttt{(sin:cos:tan)} est
DEUX fois sous une rotation 180°, donc l'unrotation doit cascader.


\begin{lstlisting}
(_,1 : +) ~ + ~ (sin : cos : tan)
\end{lstlisting}


\section{Récursion multi-fils}

Vérifie que les bundles de feedback s'imbriquent sans se croiser : l'entrée
0 de B (visible en bas après rotation 180°) reçoit la sortie 0 de A (en
haut), l'entrée 1 (au-dessus) reçoit la sortie 1, etc.


\begin{lstlisting}
A[3,3] ~ B[2,2]
\end{lstlisting}


\section{Séquentiel pur}

\begin{lstlisting}
sin : abs
\end{lstlisting}


\section{Parallèle pur}

\begin{lstlisting}
sin , cos , tan
\end{lstlisting}


\section{Split (fan-out)}

Une seule sortie distribuée vers les deux entrées de \texttt{+} :


\begin{lstlisting}
1 <: +
\end{lstlisting}


\section{Merge (fan-in)}

Quatre sorties sommées vers les deux entrées d'une boîte \texttt{(2,1)} :


\begin{lstlisting}
1, 2, 3, 4 :> +
\end{lstlisting}


\section{Primitives avec arité explicite}

\begin{lstlisting}
A[2,3] : B[3,1]
\end{lstlisting}


\section{Label entre guillemets}

\begin{lstlisting}
"my filter"[2,1]
\end{lstlisting}


\section{Tests étendus}

\subsection{Deux compteurs indépendants en parallèle}

\texttt{+\textasciitilde{}\_} est un compteur (équivalent Faust \texttt{+ \textasciitilde{} \_}). On en met deux en
parallèle, chacun alimenté par sa propre constante.


\begin{lstlisting}
1, 2 : +~_ , -~_
\end{lstlisting}


\subsection{Cascade de trois compteurs}

Compteur de compteur de compteur. Avec délais visibles.


\begin{lstlisting}
1 : +~_ : +~_ : +~_
\end{lstlisting}


\subsection{Cross-wiring : split / merge à grosse arité}

\begin{lstlisting}
(sin, cos) <: (+, *, +, *) :> (_, _)
\end{lstlisting}


\subsection{Boucle complexe : split/merge dans une rec multi-fils}

\begin{lstlisting}
((_,_) <: (_,_,_,_) :> (+,+)) ~ (mem, mem)
\end{lstlisting}


\subsection{Croisement de deux câbles}

Un grand classique BDA : on permute deux signaux en exploitant le modulo
du split. \texttt{out 0} part en position 2 (récupéré par le second \texttt{\_}), \texttt{out 1}
part en position 1 (récupéré par le premier \texttt{\_}), et les copies en
position 0 et 3 sont jetées dans les \texttt{!}.


\begin{lstlisting}
_,_ <: !,_,_,!
\end{lstlisting}


\subsection{Rec à grosses arités avec délais}

\texttt{Foo[4,5] \textasciitilde{} Bar[3,4]} — 4 entrées de Foo dont 4 consommées par feedback
(0 entrée exposée), 5 sorties dont 3 alimentent Bar, 5 exposées.


\begin{lstlisting}
Foo[4,5] ~ Bar[3,4]
\end{lstlisting}


\subsection{Passthrough plus long}

Une chaîne de 8 identités au milieu de deux blocs. Doit rester plat à
gauche (passthrough), avec des elbows propres à droite.


\begin{lstlisting}
A[3,8] : _,_,_,_,_,_,_,_ : B[8,2]
\end{lstlisting}


\subsection{Mix de tout}

\begin{lstlisting}
(1, _ : +) ~ (sin : cos) , (2, _ : *) ~ (tan : exp)
\end{lstlisting}


\section{Erreur de typage}

Le \texttt{:} exige autant de sorties à gauche que d'entrées à droite. Ici
\texttt{1} a 1 sortie mais \texttt{+} attend 2 entrées :


\begin{lstlisting}
1 : +
\end{lstlisting}

\end{document}
